% !TEX program = xelatex
\documentclass[11pt,a4paper]{article}
\usepackage{algo-preamble}

\title{\textbf{L07 \textperiodcentered{} Γραφήματα II}\\[2pt]
\large Διάσχιση --- BFS \& DFS}
\author{Algorithms Class Hub}
\date{\small Αλγόριθμοι και Πολυπλοκότητα (K17) \textperiodcentered{} ΕΚΠΑ}

\begin{document}
\maketitle

\begin{center}\itshape\small
Οι δύο θεμελιώδεις αλγόριθμοι διάσχισης γραφήματος: αναζήτηση κατά βάθος (DFS)
και αναζήτηση κατά πλάτος (BFS). Λύνουν s--t συνεκτικότητα και s--t απόσταση σε
χρόνο $O(n+m)$.
\end{center}

\section{Τι θα δούμε}

Από το προηγούμενο μάθημα ξέρουμε τι είναι ένα γράφημα και πώς το αποθηκεύουμε.
Τώρα ήρθε η ώρα για \textbf{αλγορίθμους}. Σε αυτή τη διάλεξη θα δούμε τις
\textbf{δύο θεμελιώδεις τεχνικές διάσχισης}:
\begin{enumerate}
  \item \textbf{DFS (Depth-First Search)} --- αναζήτηση κατά βάθος. Πας όσο πιο
  μακριά μπορείς, μετά γυρνάς πίσω και δοκιμάζεις άλλο δρόμο.
  \item \textbf{BFS (Breadth-First Search)} --- αναζήτηση κατά πλάτος.
  Εξερευνείς \textbf{επίπεδα} --- όλους τους γείτονες πρώτα, μετά τους γείτονες
  των γειτόνων, κ.ο.κ.
\end{enumerate}

\textbf{Και οι δύο} τρέχουν σε χρόνο $O(n + m)$ με λίστα γειτνίασης. Λύνουν την
\textbf{s--t συνεκτικότητα}. Επιπλέον το BFS υπολογίζει τη \textbf{συντομότερη
απόσταση} σε γραφήματα \textbf{χωρίς βάρη} (κάθε ακμή μετράει 1).

\begin{keypoint}
\textbf{Από τις σημαντικότερες διαλέξεις του μαθήματος.} Σχεδόν κάθε εξεταζόμενο
θέμα σε γραφήματα ξεκινάει με «τρέξε BFS» ή «τρέξε DFS» --- οι ίδιοι αλγόριθμοι ή
παραλλαγές τους είναι τα δομικά μπλοκ για τοπολογική ταξινόμηση, διμερότητα, SCC,
και πολλά άλλα που θα δούμε στα επόμενα μαθήματα.
\end{keypoint}

\section{DFS --- Αναζήτηση κατά βάθος}

\subsection{Διαισθητική ιδέα}

Φαντάσου ότι είσαι στον κόμβο $s$ ενός λαβύρινθου. Διαλέγεις \textbf{μία}
κατεύθυνση και πας. Φτάνεις σε νέο κόμβο, διαλέγεις \textbf{μία} κατεύθυνση που
δεν ξέρεις ακόμα, πας. Συνεχίζεις μέχρι να φτάσεις σε αδιέξοδο ή σε κόμβο που
έχεις ξαναδεί. Τότε \textbf{γυρνάς πίσω} στον προηγούμενο κόμβο και δοκιμάζεις
άλλη κατεύθυνση.

\textbf{Με μία λέξη:} βυθίσου όσο πιο βαθιά μπορείς, μετά γύρνα πίσω.

\subsection{Γενικός σκελετός --- η συνεκτική συνιστώσα}

Στις διαφάνειες, το «DFS» στην πρώτη του εκδοχή είναι ένας \textbf{γενικός
αλγόριθμος εύρεσης συνεκτικής συνιστώσας} --- ο \textbf{R-set algorithm}:

\begin{codebox}
\begin{verbatim}
DFS(G, s):
  R = {s}
  ενόσω υπάρχει ακμή {u, v} με u ∈ R και v ∉ R:
    R <- R ∪ {v}
  επίστρεψε R
\end{verbatim}
\end{codebox}

\textbf{Θεώρημα.} Με τον τερματισμό, το $R$ είναι η \textbf{συνεκτική συνιστώσα}
που περιέχει την $s$.

\textbf{Διαισθητικά:} ξεκινάμε από το $\{s\}$. Όσο βρίσκουμε κάποια ακμή που έχει
το ένα άκρο μέσα στο $R$ και το άλλο έξω, την «τραβάμε μέσα». Όταν δεν υπάρχει
τέτοια ακμή, το $R$ είναι ακριβώς το σύνολο των κορυφών που μπορούμε να φτάσουμε
από την $s$.

\begin{intuition}
Αυτό \textbf{δεν} είναι ο κανονικός αναδρομικός DFS, είναι μια γενική περιγραφή.
Ο πραγματικός DFS διαλέγει συγκεκριμένη σειρά (βάθος πρώτα) και υλοποιείται είτε
αναδρομικά είτε με στοίβα. Αλλά η αλήθεια είναι: \textbf{όποια σειρά κι αν
διαλέξεις} για την επόμενη ακμή, καταλήγεις στο ίδιο $R$.
\end{intuition}

\subsection{Αναδρομική υλοποίηση}

Η συνηθισμένη εκδοχή του DFS, με χρήση αναδρομής:

\begin{codebox}
\begin{verbatim}
DFS(G, s):
  R = ∅
  Για κάθε u ∈ V:
    Εξερευνημένη(u) = false
  Εξερευνημένη(s) = true
  R = {s}
  Για κάθε ακμή {s, u} στο G:
    Εάν Εξερευνημένη(u) = false:
      DFS(G, u)        // αναδρομική κλήση
\end{verbatim}
\end{codebox}

\textbf{Πώς λειτουργεί:}
\begin{enumerate}
  \item Σημαδεύουμε κάθε κορυφή ως «μη εξερευνημένη» αρχικά.
  \item Σημαδεύουμε την $s$ ως εξερευνημένη.
  \item Για \textbf{κάθε} γειτονική κορυφή της $s$ που δεν έχουμε
  ξανα-επισκεφθεί, \textbf{κατεβαίνουμε} εκεί (αναδρομική κλήση).
\end{enumerate}

\textbf{Παράδειγμα τρεξίματος} στο γράφημα του προηγούμενου μαθήματος, ξεκινώντας
από το 1:
\begin{itemize}
  \item Σημαδεύω 1. Γείτονες του 1: $\{2, 3\}$.
  \item Πάω στο 2. Γείτονες του 2: $\{1, 3, 4, 5\}$ --- το 1 εξερευνημένο.
  \item Πάω στο 3. Γείτονες του 3: $\{1, 2, 5, 7, 8\}$ --- 1, 2 εξερευνημένα.
  \item Πάω στο 5. Γείτονες του 5: $\{2, 3, 4, 6\}$ --- 2, 3 εξερευνημένα.
  \item Πάω στο 4, μετά στο 6 (αδιέξοδα) --- πίσω. Πάω στο 7, μετά στο 8.
  \item Όλη η συνιστώσα εξερευνήθηκε.
\end{itemize}
\textbf{Σειρά εξερεύνησης:} $1 \to 2 \to 3 \to 5 \to 4 \to 6 \to 7 \to 8$.

\subsection{Πολυπλοκότητα DFS: $O(n + m)$}

Με λίστα γειτνίασης: κάθε κορυφή σημαδεύεται \textbf{μία φορά} (σύνολο $O(n)$),
και κάθε ακμή εξετάζεται \textbf{δύο φορές} (μία από κάθε άκρο, σύνολο $O(m)$).
Συνολικά $O(n + m)$.

\section{BFS --- Αναζήτηση κατά πλάτος}

\subsection{Διαισθητική ιδέα}

Φαντάσου τώρα ότι είσαι \textbf{κύμα} που εκπέμπεται από την $s$. Το κύμα
εξαπλώνεται \textbf{σε όλες τις κατευθύνσεις ταυτόχρονα}. Σε χρόνο 1, φτάνεις
στους άμεσους γείτονες. Σε χρόνο 2, στους γείτονες των γειτόνων. Κ.ο.κ.

\textbf{Με μία λέξη:} εξερεύνηση \textbf{επίπεδο-επίπεδο}.

\subsection{Επίσημος ορισμός}

Από την $s$ δημιουργούμε \textbf{επίπεδα}:
\begin{itemize}
  \item $L_0 = \{s\}$.
  \item $L_1 = $ όλοι οι γείτονες του $L_0$ που δεν είναι στο $L_0$.
  \item $L_2 = $ όλες οι κορυφές με ακμή προς $L_1$, που δεν είναι σε $L_0$ ή
  $L_1$.
  \item Γενικά: $L_{i+1} = $ όλες οι κορυφές με ακμή προς $L_i$, που δεν είναι
  σε προηγούμενο επίπεδο.
\end{itemize}

\begin{keypoint}
\textbf{Θεώρημα.} Για κάθε $i$, το $L_i$ περιέχει \textbf{ακριβώς} τις κορυφές σε
\textbf{απόσταση} $i$ από την $s$ (όπου απόσταση $=$ αριθμός ακμών συντομότερης
διαδρομής).

\textbf{Διαισθητικά:} μια κορυφή σε απόσταση 3 δεν μπορεί να εμφανιστεί νωρίτερα
από το $L_3$, αλλά εμφανίζεται \textbf{εξάπαντος} στο $L_3$ (αφού έχει διαδρομή
μήκους 3). Η απόδειξη με επαγωγή στα επίπεδα είναι άμεση.

\textbf{Συνέπεια:} το BFS λύνει το πρόβλημα \textbf{s--t απόσταση} σε γραφήματα
χωρίς βάρη: η απόσταση είναι το επίπεδο όπου εμφανίζεται η $t$. Υπάρχει διαδρομή
από $s$ σε $t$ αν και μόνο αν η $t$ εμφανίζεται σε \textbf{κάποιο} επίπεδο.
\end{keypoint}

\subsection{Σημαντική ιδιότητα του BFS-δένδρου}

\begin{keypoint}
\textbf{Λήμμα.} Έστω $T$ το BFS-δένδρο που παράγεται ξεκινώντας από την $s$. Αν
$\{x, y\}$ ακμή του $G$, τότε τα \textbf{επίπεδα} των $x$ και $y$ διαφέρουν
\textbf{το πολύ κατά 1}.

\textbf{Διαισθητικά:} αν είχαμε ακμή που πηδάει 2 επίπεδα, θα έπρεπε ο γείτονας
του χαμηλότερου επιπέδου να είχε ήδη ανακαλυφθεί ένα επίπεδο νωρίτερα ---
αντίφαση.

\textbf{Συνέπεια:} όλες οι ακμές του γραφήματος είτε είναι ακμές \textbf{του
δέντρου} (μεταξύ διαδοχικών επιπέδων), είτε είναι ακμές \textbf{μέσα στο ίδιο
επίπεδο}. Καμία ακμή δεν πηδάει επίπεδα. Αυτή η ιδιότητα είναι κλειδί για τη
διμερότητα, που θα δούμε παρακάτω.
\end{keypoint}

\subsection{Ο αλγόριθμος BFS (αναλυτικά)}

\begin{codebox}
\begin{verbatim}
BFS(G, s):
  Για κάθε v ∈ V:
    Αναγνωσμένη(v) = false
    π[v] <- ∅
  Αναγνωσμένη(s) = true
  L[0] = {s}
  i <- 0           // δένδρο T = {s}
  ενόσω L[i] ≠ ∅:
    αρχικοποίηση κενής λίστας L[i+1]
    για κάθε u ∈ L[i]:
      για κάθε ακμή {u, v} ∈ E:
        εάν Αναγνωσμένη(v) = false:
          Αναγνωσμένη(v) = true
          π[v] <- u                 // T = T ∪ {u, v}
          L[i+1] <- L[i+1] ∪ {v}
    i <- i + 1
\end{verbatim}
\end{codebox}

\textbf{Τι κρατάει ο αλγόριθμος:}
\begin{itemize}
  \item \texttt{Αναγνωσμένη(v)} --- έχει επισκεφθεί η κορυφή.
  \item \texttt{π[v]} --- ο \textbf{γονέας} της $v$ στο BFS-δένδρο. Με αυτό
  μπορούμε να ανακατασκευάσουμε τη συντομότερη διαδρομή από $s$ σε κάθε κορυφή.
  \item \texttt{L[i]} --- η λίστα κορυφών στο επίπεδο $i$.
\end{itemize}

\textbf{Παραλλαγή με ουρά (queue):} αντί να φυλάμε πίνακα από λίστες, μπορούμε να
χρησιμοποιήσουμε \textbf{FIFO ουρά}. Βγάζεις την κορυφή από την αρχή, βάζεις τους
μη-εξερευνημένους γείτονες στο τέλος. Είναι μαθηματικά \textbf{ισοδύναμο}.

\subsection{Παράδειγμα στο γράφημα του προηγούμενου μαθήματος}

Ξεκινώντας από $s = 1$:
\begin{itemize}
  \item $L_0 = \{1\}$.
  \item Γείτονες του 1: $\{2, 3\}$. Άρα $L_1 = \{2, 3\}$.
  \item Γείτονες του 2: $\{1, 3, 4, 5\}$ (νέοι: 4, 5)· γείτονες του 3:
  $\{1, 2, 5, 7, 8\}$ (νέοι: 7, 8). Άρα $L_2 = \{4, 5, 7, 8\}$.
  \item Γείτονες του 5: $\{2, 3, 4, 6\}$ (νέο: 6). Άρα $L_3 = \{6\}$.
\end{itemize}

\textbf{Αποστάσεις από το 1:}
\begin{center}
\begin{tabular}{@{}l c c c c c c c c@{}}
\toprule
Κορυφή   & 1 & 2 & 3 & 4 & 5 & 6 & 7 & 8 \\
\midrule
Απόσταση & 0 & 1 & 1 & 2 & 2 & 3 & 2 & 2 \\
\bottomrule
\end{tabular}
\end{center}

\subsection{Ανάλυση πολυπλοκότητας}

Δείχνουμε ότι ο BFS τρέχει σε $O(n + m)$ με λίστα γειτνίασης.

\textbf{Αφελής φραγμός} $O(n^2)$: το πολύ $n$ λίστες $L[i]$· κάθε κορυφή ανήκει
σε το πολύ μία· για κάθε κορυφή $u$ ο βρόχος «για κάθε ακμή $\{u, v\}$» κάνει το
πολύ $n$ επαναλήψεις· $\Theta(1)$ ανά επανάληψη $\to O(n^2)$ σύνολο.

\textbf{Σφιχτότερος φραγμός} $O(n + m)$: όταν εξετάζουμε την κορυφή $u$,
σαρώνουμε \textbf{ακριβώς $\deg(u)$} προσκείμενες ακμές (όχι $n$). Σύνολο
εξετάσεων:
\[
\sum_{v \in V} (1 + \deg(v)) = n + \sum_{v} \deg(v) = n + 2m
\]
(από το handshaking λήμμα του προηγούμενου μαθήματος). Άρα \textbf{$O(n + m)$}.

\begin{keypoint}
\textbf{Σημαντική παρατήρηση:} η ανάλυση $O(n^2)$ είναι \textbf{σωστή}
ασυμπτωτικά αλλά \textbf{όχι σφιχτή}. Για αραιά γραφήματα ($m = O(n)$) ο σφιχτός
φραγμός $O(n+m) = O(n)$ είναι \textbf{πολύ καλύτερος}. Πάντα ψάχνε σφιχτό φραγμό
σε ανάλυση γραφημάτων.
\end{keypoint}

\section{Σύγκριση: BFS vs DFS}

\begin{center}
\begin{tabular}{@{}l p{4.7cm} p{4.7cm}@{}}
\toprule
 & \textbf{BFS} & \textbf{DFS} \\
\midrule
Στρατηγική        & επίπεδο-επίπεδο & βάθος πρώτα \\
Δομή δεδομένων    & ουρά (FIFO)     & στοίβα (LIFO) ή αναδρομή \\
Λύνει             & s--t απόσταση (χωρίς βάρη) & s--t συνεκτικότητα \\
Πολυπλοκότητα     & $O(n + m)$      & $O(n + m)$ \\
Δένδρο που παράγει& ίδιο επίπεδο $=$ ίδια απόσταση & μακρές αλυσίδες \\
Εφαρμογές         & shortest paths (unweighted), bipartite test, web crawler & τοπολογική ταξινόμηση, SCC, εύρεση κύκλων \\
\bottomrule
\end{tabular}
\end{center}

\begin{intuition}
\textbf{Πότε ποιο;}
\begin{itemize}
  \item Θέλω \textbf{απόσταση}: BFS.
  \item Θέλω απλά «\textbf{υπάρχει διαδρομή;}»: είτε δουλεύει, αλλά DFS συχνά πιο
  φυσικό.
  \item Θέλω \textbf{εύρεση κύκλου} ή \textbf{τοπολογική ταξινόμηση}: DFS.
  \item Θέλω να φτάσω \textbf{όσο πιο μακριά μπορώ}: DFS.
  \item Θέλω «\textbf{ποιες κορυφές απέχουν ακριβώς $k$ βήματα;}»: BFS.
\end{itemize}
\end{intuition}

\section{Σχέση BFS / DFS με τα προβλήματα του προηγούμενου μαθήματος}

\begin{center}
\begin{tabular}{@{}p{4.6cm} p{6.2cm} l@{}}
\toprule
Πρόβλημα & Λύση & Πολυπλοκότητα \\
\midrule
s--t συνεκτικότητα & DFS ή BFS --- δες αν το $t$ ανήκει στο $R$ ή σε κάποιο
$L_i$ & $O(n + m)$ \\
s--t απόσταση (χωρίς βάρη) & BFS --- απόσταση $=$ επίπεδο & $O(n + m)$ \\
Εύρεση \textbf{όλων} των συνεκτικών συνιστωσών & Επανέλαβε DFS/BFS από κάθε
ανεξερεύνητη κορυφή & $O(n + m)$ \\
\bottomrule
\end{tabular}
\end{center}

\begin{recapbox}
\textbf{Τι κρατάμε από αυτή τη διάλεξη}
\begin{enumerate}
  \item \textbf{DFS} --- βυθίζεσαι όσο μπορείς, μετά γυρνάς. Αναδρομική
  υλοποίηση. $O(n+m)$.
  \item \textbf{BFS} --- εξερευνάς επίπεδα. $L_0 = \{s\}$, $L_{i+1} = $ νέοι
  γείτονες του $L_i$. $O(n+m)$.
  \item \textbf{Θεώρημα BFS:} $L_i = $ όλες οι κορυφές σε απόσταση $i$.
  \textbf{Συνέπεια:} BFS λύνει s--t απόσταση σε γραφήματα χωρίς βάρη.
  \item \textbf{Ιδιότητα BFS-δέντρου:} καμία ακμή δεν πηδάει επίπεδα. Κλειδί για
  τη διμερότητα.
  \item \textbf{Σύγκριση:} BFS για απόσταση, DFS για βάθος/κύκλους. Και τα δύο
  σε $O(n+m)$ με λίστα γειτνίασης.
  \item \textbf{Σφιχτή ανάλυση:} μην αρκείσαι σε $O(n^2)$ --- με handshaking
  παίρνεις $O(n+m)$, που για αραιά γραφήματα είναι πολύ καλύτερο.
\end{enumerate}
\end{recapbox}

\lecturefooter
\end{document}
